\documentclass[
	12pt,						 % Schriftgröße
	DIV=10,						 % Satzspiegelgröße Seitenränder
	german,						 % für Umlaute, Silbentrennung etc.
	paper=a4,					 % Papierformat
	oneside, 
	titlepage,					 % es wird eine Titelseite verwendet
	parskip=half,				 % Abstand zwischen Absätzen (halbe Zeile)
	headings=normal,			 % Größe der Überschriften verkleinern
	listof=totoc,				 % Verzeichnisse im Inhaltsverzeichnis aufführen
	bibliography=totoc,			 % Literaturverzeichnis im Inhaltsverzeichnis aufführen
	numbers=noenddot,			 % Kapitelnummern ohne Punkt	
	enabledeprecatedfontcommands % Alte Befehle zur unterstüzung von \listofsymbols	
]{scrreprt}

\makeatletter	
	\setlength\@fptop{0\p@}  % Entfernt Lücken am Seitenbeginn
\makeatother

\newcommand{\titel}{Entwurf einer Datenbank für ein Verwaltungssystem von Lebensmitteln in einem Kühlschrank}
\newcommand{\art}{Ausarbeitung}
\newcommand{\autor}{Claas Meints}
\newcommand{\matrikel}{172489}
\newcommand{\semester}{6}
\newcommand{\studienbereich}{Elektrotechnik/Mechatronik}
\newcommand{\studiengang}{Elektrotechnik}
\newcommand{\erstgutachter}{Prof. Dr. Herwig Henseler}
\newcommand{\jahr}{2022}
\newcommand{\comment}[1]{}
\newcommand{\TODO}{}
\newcommand\fig[2]{
	\begin{figure}[ht]
		\centering
		\includegraphics[width=0.8\textwidth]{figures/#1}
		\caption{#2} 
		\label{fig:#1}
	\end{figure}
}
% Zur Richtigen Verwendung von Einheiten
\usepackage[
locale=DE, 
%per=symbol-or-fraction
]{siunitx}

\sisetup{
	mode = text,
	detect-family,
	detect-weight,  
	exponent-product = \cdot,
	output-decimal-marker={\text{.}},
	%math-rm=\mathsf,
	%text-rm=\sffamily,
	range-phrase = { -- },
	per-mode=symbol,
	per-symbol=/
}
\DeclareSIUnit{\var}{var}

% Anpassung des Seitenlayouts -------------------------------------------------------------% 	siehe Seitenstil.tex
% -------------------------------------------------------------
\usepackage[headsepline,automark]{scrlayer-scrpage}

\usepackage{listings}
\usepackage{color}
% C++ color scheme
\definecolor{backgroundcolor}{RGB}{237, 241, 237}
\definecolor{commentcolor}{RGB}{0,128,0}
\definecolor{keywordcolor}{RGB}{136, 134, 58}
\definecolor{identifiercolor}{RGB}{24, 79, 110}
\definecolor{numbercolor}{RGB}{128,128,128}
\definecolor{preprocessorcolor}{RGB}{128,128,128}
\definecolor{stringcolor}{RGB}{163,21,21}

\lstdefinestyle{C++}{
	backgroundcolor=\color{backgroundcolor},   
    commentstyle=\color{commentcolor},
	directivestyle=\color{preprocessorcolor},
    keywordstyle=\color{keywordcolor},
	identifierstyle=\color{identifiercolor},
    numberstyle=\tiny\color{numbercolor},
    stringstyle=\color{stringcolor},
    basicstyle=\ttfamily\footnotesize,
    breakatwhitespace=false,         
    breaklines=true,                 
    captionpos=b,                    
    keepspaces=true,                 
    numbers=left,                    
    numbersep=5pt,                  
    showspaces=false,                
    showstringspaces=false,
    showtabs=false,                  
    tabsize=2
}
\lstset{language=C++,escapechar=@,style=C++}
\usepackage{booktabs}
\usepackage{eurosym}
\usepackage{pifont}% http://ctan.org/pkg/pifont

\renewcommand{\lstlistingname}{Code} % Damit Codebeispiele mit Code betitelt werden.


% Anpassung an Landessprache -------------------------------------------------------------% 	Verwendet globale Option german siehe \documentclass
% -------------------------------------------------------------
\usepackage[ngerman]{babel}
\usepackage{scrhack}
\usepackage[utf8]{inputenc}
\usepackage{amsmath,amssymb,units}
\usepackage{wrapfig,caption, placeins}
\usepackage{multicol}
\captionsetup{
	format = plain,
	justification = centering,
	labelsep = newline,
	singlelinecheck = false,
	labelfont = bf,
	font = small
}
\usepackage{mathptmx, times, helvet, courier} % Pakete für Schriftarten
%\usepackage{mathptmx,charter,helvet,courier}
\usepackage[printonlyused, smaller]{acronym}
\usepackage{wrapfig}
\usepackage{pgfplots}
% Umlaute -------------------------------------------------------------% 		Umlaute/Sonderzeichen wie äöüß direkt im Quelltext verwenden (CodePage).
%		Erlaubt automatische Trennung von Worten mit Umlauten.
% -------------------------------------------------------------
\usepackage[utf8]{inputenc}
\usepackage[T1]{fontenc}
\usepackage{ae}			 	% "schöneres" ä
\usepackage{textcomp} 		% Euro-Zeichen etc.
\usepackage[german=quotes]{csquotes}


% Verwendung von Blindtext zur Layout gestaltung -------------------------------------------------------------
\usepackage{blindtext}

% Grafiken -------------------------------------------------------------
% 		Einbinden von Grafiken [draft oder final]
% 		Option [draft] bindet Bilder nicht ein - auch globale Option
% -------------------------------------------------------------
\usepackage{graphicx}
%\graphicspath{{Grafiken/}} 	% Dort liegen die Bilder des Dokuments


% Befehle aus AMSTeX für mathematische Symbole z.B. \boldsymbol \mathbb ----
\usepackage{amsmath,amsfonts}
\usepackage{tensor}

% Für Index-Ausgabe; \printindex -------------------------------------------------------------
\usepackage{makeidx}

% Einfache Definition der Zeilenabstände und Seitenränder -------------------------------------------------------------
\usepackage{setspace}
\usepackage{geometry}


%% Zum Umfließen von Bildern -------------------------------------------------------------
\usepackage{floatflt}


% Lange URLs umbrechen etc. -------------------------------------------------------------
\usepackage{url}

% für lange Tabellen
\usepackage{longtable}
\usepackage{multirow}
\usepackage{rotating}
\usepackage{diagbox}
\usepackage{array}
\newcolumntype{M}[1]{>{\centering\arraybackslash}m{#1}}
\usepackage{ragged2e}
\usepackage{lscape}
\usepackage{colortbl} %farbige hinterlegung


% Spaltendefinition rechtsbündig mit definierter Breite auf dem Deckblatt 
%------------------------------------------------------------
\newcolumntype{d}[1]{>{\raggedleft\hspace{0pt}}p{#1}}

% Formatierung von Listen ändern
\usepackage{paralist}
% Standardeinstellungen:
\setdefaultleftmargin{2.5em}{2.2em}{1.87em}{1.7em}{1em}{1em}

%Zur Verwendung von BibLaTex
\usepackage[
	style=alphabetic
]{biblatex}
\addbibresource{bibliografie.bib}
\DefineBibliographyStrings{ngerman}{ 
	andothers = {{et\,al\adddot}},             
}

\usepackage{xcolor}
\usepackage{caption}
\usepackage{float}
\usepackage{subfig}
\usepackage{xspace}
\usepackage{scrhack}

\usepackage[acronyms,toc]{glossaries} %Glossar Achtung Perl.exe muss installiert werden
\usepackage{glossaries-extra}
\usepackage{glossary-mcols}

% pdf-Optionen  -------------------------------------------------------------
\usepackage{pgfplots} % Für Kompatibilät unter verschiedenen OS
\pgfplotsset{compat=1.7}

\pdfminorversion=7 % Verwende pdf Version 1.7

\usepackage[
bookmarks,
bookmarksopen=true,
bookmarksnumbered=true,
colorlinks=true,
%linkcolor=red, % einfache interne Verknüpfungen
%anchorcolor=black,% Ankertext
%citecolor=blue, % Verweise auf Literaturverzeichniseinträge im Text
%filecolor=magenta, % Verknüpfungen, die lokale Dateien öffnen
%menucolor=red, % Acrobat-Menüpunkte
%urlcolor=cyan, 
% für die Druckversion können die Farben ausgeschaltet werden:
linkcolor=black, % einfache interne Verknüpfungen
anchorcolor=black,% Ankertext
citecolor=black, % Verweise auf Literaturverzeichniseinträge im Text
filecolor=black, % Verknüpfungen, die lokale Dateien öffnen
menucolor=black, % Acrobat-Menüpunkte
urlcolor=black, 
%backref,
%pagebackref,
plainpages=false,% zur korrekten Erstellung der Bookmarks
pdfpagelabels,% zur korrekten Erstellung der Bookmarks
hypertexnames=false,% zur korrekten Erstellung der Bookmarks
linktoc=all%Sowohl Seitenzahl als auch Text als Link
]{hyperref}
\usepackage{bookmark}
\usepackage[all]{hypcap} % jump to top of figure
\usepackage{pdfpages}
% Zeilenabstand ------------------------------------------------------------
\onehalfspacing

% Schriftart ---------------------------------------------------------------
% \renewcommand{\familydefault}{\sfdefault} % Schriftart

% Seitenränder -------------------------------------------------------------
\geometry{paper=a4paper,left=25mm,right=25mm,top=30mm,bottom=40mm}

% Kopf- und Fußzeilen ------------------------------------------------------
\pagestyle{scrheadings}

% Kopf- und Fußzeile auch auf Kapitelanfangsseiten -------------------------
\renewcommand*{\chapterpagestyle}{scrheadings}

% Schriftform der Kopfzeile ------------------------------------------------
\renewcommand{\headfont}{\normalfont}

% Kopfzeile ----------------------------------------------------------------
\ihead{
	\parbox{11.8cm}{\normalsize{ \ }}\\
	%\small{\untertitel}\\[2ex]
	\textit{ \headmark}
}
\chead{
		%\includegraphics[height=0.8cm]{Grafiken/EWENETZ_Logo.png}
}
\ohead{	
		\includegraphics[height=1.4cm]{figures/Logo_PHWT.png}	
}
\setlength{\headheight}{21mm} % Höhe der Kopfzeile


% Fußzeile -----------------------------------------------------------------
\ifoot{
	%\includegraphics[height=1.5cm]{Grafiken/Logo_PHWT.png}
}
\cfoot{}
\ofoot{\pagemark}
\setlength{\footheight}{15mm}


% erzeugt ein wenig mehr Platz hinter einem Punkt --------------------------
%\frenchspacing 

% Schusterjungen und Hurenkinder vermeiden
\clubpenalty = 10000
\widowpenalty = 10000 
\displaywidowpenalty = 10000


%% Fußnoten fortlaufend durchnummerieren ------------------------------------
%\counterwithout{footnote}{chapter}
\setkomafont{section}{\large} 



\newglossary[llg]{glossary}{llo}{lls}{Glossar}
\makeglossaries

\begin{document}
	\thispagestyle{plain}
\begin{titlepage}
    \begin{center}
        \huge{\textbf{\art}}\\[1.5ex]
        \LARGE{\titel}\\[2ex]
        \large{Private Hochschule für Wirtschaft und Technik Vechta/Diepholz}
        \begin{tabular}{r l}
            vorgelegt von: & \quad \autor \\
            Semester: & \quad \semester \\
        \end{tabular}
    \end{center}
    \setcounter{tocdepth}{0}\let\clearpage\relax\tableofcontents
    {\centering\textcopyright\ 2022\\[1.5ex]}
    \singlespacing\small\noindent Dieses Werk einschließlich seiner Teile ist \textbf{urheberrechtlich geschützt}. Jede Verwertung außerhalb der engen Grenzen des Urheberrechtgesetzes ist ohne Zustimmung des Autors unzulässig und strafbar. Das gilt insbesondere für Vervielfältigungen, Übersetzungen, Mikroverfilmungen sowie die Einspeicherung und Verarbeitung in elektronischen Systemen. 
\end{titlepage}

	
	\pagenumbering{Roman}
	\setcounter{tocdepth}{2} % Tiefe des Inhaltsverzeichnisses definiert
	\tableofcontents

    \chapter{Einleitung}\label{ch:Einleitung}
\pagenumbering{arabic}
\section{Motivation}\label{sec:Motivation}
- IoT ist ein großes Thema\\
- Kennzahlen für IoT\\
\section{Ziel der Arbeit}\label{sec:Ziel der Arbeit}
- Ausführlicher Vergleich von IoT-Plattformen mit Fokus auf die Entwicklung von IoT-Lösungen vom Prototyp bis zum skalierten Geschäftsmodell\\
- Auswahl geeigneter Kriterien anhand eines Referenzmodells mit Rückblick auf Prototypentwicklung und Skalierung\\
\section{Aufbau der Arbeit}\label{sec:Aufbau der Arbeit}

    \chapter{Stand der Technik}\label{ch:Stand der Technik}
\section{Komponenten einer IoT-Lösung}\label{sec:Komponenten einer IoT-Lösung}
% [Guth2018] [10.1007/978-3-030-24311-1_8]\\
\subsection{Geräte}\label{subsec:Geräte}
- Sensoren und Aktoren\\
- Treiber für Sensoren und Aktoren\\
- Geräte SDKs\\
- RTOS\\
- Edge Computing\\
\subsection{Gateway}\label{subsec:Gateway}
- Software und Schnittstellen für Gerätekommunikation\\
- Gateway SDKs\\
- OS\\
- Schnittstellen zu Cloud-Plattformen (MQTT, REST, etc.)\\
- Fog Computing\\
\subsection{IoT-Software-Plattform}\label{subsec:IoT-Software-Plattform}
- Verwaltet Geräte und Gateways\\
- Stellt Schnittstellen für externe Applikationen und Services bereit\\
- bietet Services für die Verarbeitung von Daten\\
- wird in dieser Arbeit als IoT-Plattform bezeichnet\\
\subsection{Anwendung}\label{subsec:Anwendung}
- Anwendungen / Applikationen / Services auf Basis der IoT-Plattform und der Daten der Geräte\\
\section{Referenzarchitektur von IoT-Software-Plattformen}\label{sec:Referenzarchitektur von IoT-Software-Plattformen}
\subsection{Kommunikation}\label{subsec:Kommunikation}
- MQTT Broker\\
- (MQTT-)WebSockets\\
- HTTP (REST)\\
\subsection{Digitaler Zwilling}\label{subsec:Digitaler Zwilling}
- Device Registry / Twin\\
- Device Metadata and Data\\
- Remote Actions\\
- Device Updates\\
\subsection{IoT-Services}\label{subsec:IoT-Services}
% [Lempert2021] [Microsoft_2021] [9825678]\\
- Develop Insights (Cold, Warm and Hot Path)\\
- Databases (Structured, Unstructured and Time Series)\\
- Descriptive Analytics\\
- Diagnostic Analytics\\
- Real-Time Analytics (Stream Processing Hot Path)\\
- Predictive / Prescriptive Analytics (Machine Learning Cold Path)\\
\subsection{Integration in Drittsysteme}\label{subsec:Integration in Drittsysteme}
- Adapter zu Enterprise-Systemen\\
- Cloud to Cloud Integration\\
- Apis?\\
\subsection{Orchestrierung von Services}\label{subsec:Orchestrierung von Services}
- Message Broker / Event Bus\\
- Rules Engine\\
- Service-Orchestrierung\\
- Load Balancing?\\
\subsection{Management}\label{subsec:Management}
- Device Management -> Device Twin\\
- Configuration / State Management -> Device Twin\\
- Errors and Reports\\
\subsection{Security}\label{subsec:Security}
- Authentication / Authorization\\
- Identity Management\\
\section{Aktuelle Implementierungen von IoT-Software-Plattformen}\label{sec:Aktuelle Implementierungen von IoT-Software-Plattformen}
\subsection{Hyperscaler-Cloud-Anbieter}\label{subsec:Hyperscaler-Cloud-Anbieter}
- Bindung an den Cloud-Anbieter?\\
- Regional verteilbar -> geringere Latenz?\\
\subsection{(Open-Source IoT-Plattformen)}\label{subsec:(Open-Source IoT-Plattformen)}
- Cloud agnostisch?\\
- On-Premise möglich?\\
\subsection{(IoT-Spezifische Plattformen)}\label{subsec:(IoT-Spezifische Plattformen)}
% ^ möglicherweise zusammenfassen ^\\
\subsection{Low-Code-Plattformen}\label{subsec:Low-Code-Plattformen}
- schnellere Entwicklung von IoT-Lösungen?\\
- höhere Flexibilität?\\
- hohe Abhängigkeit von der Plattform?\\
- hohe Kosten?\\
\section{IoT Geschäftsmodellentwicklung}\label{sec:IoT Geschäftsmodellentwicklung}
\subsection{Prototypentwicklung}\label{subsec:Prototypentwicklung}
\subsection{Skalierung}\label{subsec:Skalierung}

    \chapter{Theoretische Grundlagen}\label{ch:Theoretische Grundlagen}
\section{(Übertragungstechnologien)}\label{sec:(Übertragungstechnologien)}
\subsection{(WLAN)}\label{subsec:(WLAN)}
\subsection{(Mobilfunk)}\label{subsec:(Mobilfunk)}
\subsection{(LoRaWAN)}\label{subsec:(LoRaWAN)}
\section{(Übertragungsprotokolle und Schnittstellen)}\label{sec:(Übertragungsprotokolle und Schnittstellen)}
\subsection{(MQTT)}\label{subsec:(MQTT)}
- Publish / Subscribe\\
- MQTT WebSockets\\
\subsection{(HTTP)}\label{subsec:(HTTP)}
- REST\\

    \chapter{Anforderungsdefiniton (An die Arbeit)}\label{ch:Anforderungsdefiniton (An die Arbeit)}
\section{Entwicklung von Kriterien zur Auswahl einer IoT-Software-Plattform}\label{sec:Entwicklung von Kriterien zur Auswahl einer IoT-Software-Plattform}
\section{Vergleich der technischen Funktionalitäten von IoT-Software-Plattformen}\label{sec:Vergleich der technischen Funktionalitäten von IoT-Software-Plattformen}
\section{Vergleich nicht-funktionaler Anforderungen an IoT-Software-Plattformen}\label{sec:Vergleich nicht-funktionaler Anforderungen an IoT-Software-Plattformen}
\section{Betrachtung von der Prototyp-Phase bis zur skalierten Lösung}\label{sec:Betrachtung von der Prototyp-Phase bis zur skalierten Lösung}
\section{(Evaluierung anhand eines Beispielprojekts)}\label{sec:(Evaluierung anhand eines Beispielprojekts)}
\section{Abgrenzung der Arbeit}\label{sec:Abgrenzung der Arbeit}

    \input{chapters/5 Systemanalyse (Grundlage für Kriterien).tex}
    \chapter{Konzeption (Kriterien für Prototy und Skalierte Lösung)}\label{ch:Konzeption (Kriterien für Prototy und Skalierte Lösung)}
\section{Funktionale Anforderungen}\label{sec:Funktionale Anforderungen}
\subsection{Verfügbarkeit}\label{subsec:Verfügbarkeit}
- Redundancy\\
- Fault Resilience\\
- Regional Distribution\\
- Self Healing\\
\subsection{Performanz}\label{subsec:Performanz}
- Latency\\
- Data Processing / Analytics\\
\subsection{Unterstützte Funktionen und Schnittstellen}\label{subsec:Unterstützte Funktionen und Schnittstellen}
- Communication Technologies\\
- Open Adapters / APIs\\
- Implementation of optional features\\
\subsection{Sicherheit}\label{subsec:Sicherheit}
- Support for Security-Hardware\\
\subsection{Skalierbarkeit}\label{subsec:Skalierbarkeit}
\section{Nicht-funktionale Anforderungen}\label{sec:Nicht-funktionale Anforderungen}
\subsection{Kosten und Lizenzierung}\label{subsec:Kosten und Lizenzierung}
- Kosten pro Gerät\\
- Kosten per Nutzung\\
- Lizenzkosten\\
- Open Source?\\
\subsection{Rechtliche Anforderungen}\label{subsec:Rechtliche Anforderungen}
- Datenschutz\\
- Datentransfer (EU-US Privacy Shield)\\
- Kritis?\\
\subsection{Entwicklungsaufwand und -geschwindigkeit}\label{subsec:Entwicklungsaufwand und -geschwindigkeit}
- Welche Vorkenntnisse sind nötig?\\
- Wie viel Aufwand wird benötigt um einen einfachen Prototypen zu entwickeln?\\
- Wie leicht lassen sich neue Funktionen implementieren?\\
- Wie leicht lässt sich die Lösung anpassen oder migrieren?\\
- Developer Experience\\
\section{Bewertung der Anforderungen}\label{sec:Bewertung der Anforderungen}
\subsection{Kritische Anforderungen}\label{subsec:Kritische Anforderungen}
- zwingend notwendig für die Entwicklung einer IoT-Lösung\\
- gibt es zusammengesetzte kritische Anforderungen?\\
- welche Zustände sind für Projekte unakzeptabel? (in der Prototypen- und Skalierungsphase)\\
\subsection{Gewichtung der Anforderungen}\label{subsec:Gewichtung der Anforderungen}
- Wie wichtig sind die Anforderungen für das Projekt?\\
- Ändern sich Gewichtungen im Laufe des Projekts?\\

    \chapter{Implementierung (Abgleich der Kriterien)}\label{ch:Implementierung (Abgleich der Kriterien)}
\section{Vorauswahl von IoT-Software-Plattformen}\label{sec:Vorauswahl von IoT-Software-Plattformen}
- Können Plattformen oder Plattformtypen ausgeschlossen werden?\\
- Welches sind die wichtigsten und relevantesten IoT-Plattformen?\\
- Welche IoT-Plattformen sind für die EWE besonders interessant?\\
- Welche IoT-Plattformen sind für die Entwicklung  digitaler Geschäftsmodelle besonders interessant?\\
\section{Analyse der IoT-Plattformen}\label{sec:Analyse der IoT-Plattformen}
\subsection{AWS Core IoT}\label{subsec:AWS Core IoT}
\subsection{Azure IoT}\label{subsec:Azure IoT}
\subsection{Google Cloud IoT}\label{subsec:Google Cloud IoT}
\subsection{(Digimondo nIoTx)}\label{subsec:(Digimondo nIoTx)}
\subsection{(Zenner IoT)}\label{subsec:(Zenner IoT)}
\section{Vergleich der IoT-Plattformen}\label{sec:Vergleich der IoT-Plattformen}

    \input{chapters/8 Umsetzung eines exemplarischen IoT-Geschäftsmodells.tex}
    \chapter{Ergebnis und Diskussion}\label{ch:Ergebnis und Diskussion}
\section{Anwendungsspezifische Bewertung von IoT-Plattformen}\label{sec:Anwendungsspezifische Bewertung von IoT-Plattformen}
- (Fortsetzung 7.3)\\
- Welche Plattformen eignen sich für welche Anwendungsfälle?\\
- Welche Kriterien sind besonders wichtig für welche Anwendungsfälle?\\
\section{Bewertung des Auswahlprozesses anhand des Kriterienkatalogs}\label{sec:Bewertung des Auswahlprozesses anhand des Kriterienkatalogs}
- (zu 7.2/7.3)\\
- Ist der Kriterienkatalog ausreichend?\\
- Welche Schwierigkeiten ergeben siche bei der Anwendung auf die IoT-Plattformen?\\
\section{Ergebnisse der exemplarischen Umsetzung des Auswahlprozesses}\label{sec:Ergebnisse der exemplarischen Umsetzung des Auswahlprozesses}
-(zu 8)\\

    \chapter{Zusammenfassung}\label{ch:Zusammenfassung}

    \chapter{Ausblick}\label{ch:Ausblick}

 
\end{document} 